\subsection{উদাহরণ: উপস্থিতির উপাত্ত থেকে সিদ্ধান্ত}
এক সপ্তাহে একটি শ্রেণির উপস্থিত শিক্ষার্থীর সংখ্যা যথাক্রমে 42, 45, 40, 46 ও 47। এসব মান আলাদা দিনের উপাত্ত। মোট উপস্থিতি $42+45+40+46+47=220$ এবং দৈনিক গড় উপস্থিতি $220/5=44$। শ্রেণিতে মোট 50 জন থাকলে গড় উপস্থিতির হার $44/50\times100=88\%$। এখন ``এই সপ্তাহে গড়ে ৮৮ শতাংশ শিক্ষার্থী উপস্থিত ছিল'' একটি ব্যবহারযোগ্য তথ্য।

শিক্ষক শুধু গড় দেখলে কোন দিন উপস্থিতি সবচেয়ে কম ছিল তা সরাসরি নাও বুঝতে পারেন। মূল উপাত্ত থেকে দেখা যায় তৃতীয় দিনে উপস্থিতি 40। অর্থাৎ একই উপাত্ত থেকে বিভিন্ন উদ্দেশ্যের তথ্য তৈরি হয়। তাই তথ্য তৈরি করার সময় কোন প্রশ্নের উত্তর প্রয়োজন, সেটি নির্ধারণ করতে হয়।
\begin{flushleft}\small
\begin{tabularx}{\linewidth}{|p{33mm}|X|}\hline
\chSixHeader প্রশ্ন & উপাত্ত থেকে পাওয়া তথ্য\\\hline
মোট উপস্থিতি কত? & পাঁচ দিনের উপস্থিতি-গণনার যোগফল 220\\\hline
গড়ে কতজন এসেছে? & প্রতিদিন গড়ে 44 জন\\\hline
কোন দিন সবচেয়ে কম? & তৃতীয় দিন, 40 জন\\\hline
সর্বোচ্চ ও সর্বনিম্নের পার্থক্য? & $47-40=7$ জন\\\hline
\end{tabularx}
\end{flushleft}
এখানে মোট উপস্থিতি 220 মানে 220 জন পৃথক শিক্ষার্থী নয়; একই শিক্ষার্থী কয়েক দিন এসেছে। তাই গণনার ফলের সঙ্গে অর্থ ও একক উল্লেখ করাও তথ্যের নির্ভুল উপস্থাপনের অংশ।

